\section{Ordinary reflections and covariance of individual phases}
\label{ex-reflection-transport}

Reflected conjugation is built into the Hilbert inner product. Ordinary
reflection, which acts complex-linearly on labels, is a different
operation. Proving that it fixes each compressed observable is the extra
step needed to obtain orthogonal covariance in each character separately.
We denote ordinary time reflection by \(\rho\); thus \(\theta\)
also conjugates labels and reverses the word, whereas \(\rho\) does
neither. Permutation symmetry makes the reversal immaterial to \(M\).

\begin{lemma}[Positivity from bounded real approximations]
\label{ex-replicated-positivity}
Suppose real numbers \(b_t(i,j)\), for integers \(i<j\), satisfy
\(b_t(i,j)\to a\) as \(t\to\infty\), with the same limit \(a\)
for every pair. Suppose there are \(\epsilon_k\ge0\) tending to zero
and finite \(R_k\ge0\) with the following property. For every \(k\)
and finite \(m\), for all sufficiently large \(t\) there exist
simultaneous real random variables \(X_1,\ldots,X_m\), possibly on
cutoff-dependent probability spaces, such that eventually in the ordered
cutoff limit
\[
 |X_i|\le R_k,\qquad
 \mathbb E(X_iX_j)\le b_{\mathrm{cut},t}(i,j)+\epsilon_k
 \quad(i<j),
\]
where \(b_{\mathrm{cut},t}(i,j)\to b_t(i,j)\) in that cutoff limit.
Then \(a\ge0\).
\end{lemma}
\begin{proof}
The variables are bounded and measurable, so their products are integrable.
Nonnegativity of a square gives
\[
 0\le\mathbb E\left(\sum_{i=1}^mX_i\right)^2
 \le mR_k^2+2\sum_{i<j}
       \bigl(b_{\mathrm{cut},t}(i,j)+\epsilon_k\bigr).
\]
For fixed \(k,m,t\), all eventual conditions can be imposed together:
only finitely many variables and pairs occur. Remove the cutoff in the
specified order, then send \(t\to\infty\), obtaining
\(0\le mR_k^2+m(m-1)(a+\epsilon_k)\).
For \(m\ge2\), divide by \(m(m-1)\) and send \(m\to\infty\).
This yields \(a+\epsilon_k\ge0\). Now send \(k\to\infty\).
The order of these limits matters: \(R_k\) may diverge as \(k\) grows.
\end{proof}

The lattice construction applies the lemma to a real-labelled word
\(w\) at an internally generic grid configuration, minus its ordinary reflection. Translate several copies along
one vector with strictly positive coordinates, spacing the \(i\)-th
copy by \(it\) times that vector. For fixed distinct copies and large
\(t\), their combined configurations are generic in every coordinate.
The covariance to be approximated is the four-term expression
\[
 S(w_iw_j)-S(w_i(\rho w)_j)
 -S((\rho w)_iw_j)+S((\rho w)_i(\rho w)_j),
\]
or its real part. By Lemma~\ref{ex-cluster-train}, its large-separation
limit is
\begin{equation}\label{ex-reflection-square}
 a=\operatorname{Re}\langle\Omega,
       (M(w)-M(\rho w))^2\Omega\rangle.
\end{equation}
The input required here is the existence of the bounded real variables
in Lemma~\ref{ex-replicated-positivity}, with errors independent of the
number and offsets of copies. Ordinary positivity of a Wilson
probability measure supplies the nonnegative square once these
simultaneous approximations are available; it does not supply the
approximations or their uniform error bounds.

\begin{proposition}[Ordinary reflection fixes compressed words]
\label{ex-ordinary-reflection}
Assume the bounded real approximation property just described holds on
a dense set of internally generic grid configurations for each tuple of real basis labels.
Assume the norm-regularity hypotheses of Section~\ref{ex-slab-regularity}
and that each raw label space has a basis fixed by conjugation.
Then \(M(\rho w)=M(w)\) for every distinct-point word.
\end{proposition}
\begin{proof}
First take a real basis-labelled word at an admitted configuration.
Since its labels are fixed by conjugation, permutation symmetry and the
adjoint identity give \(M(\rho w)=M(\theta w)=M(w)^*\).
Thus \(D=M(w)-M(\rho w)\) satisfies \(D^*=-D\). The previous lemma
and \eqref{ex-reflection-square} imply
\[
 0\le\operatorname{Re}\langle\Omega,D^2\Omega\rangle
 =-\|D\Omega\|^2.
\]
Hence \(D\Omega=0\), and separation of \(\Omega\) for
\(\mathcal A\) gives \(D=0\). Norm continuity of compressed words
and their reflected words extends the identity from the dense admitted
configurations to all distinct configurations for these basis labels.
Both \(M\) and ordinary reflection are complex-multilinear in the
labels. Expanding an arbitrary tuple in the real basis therefore extends
the equality to every tuple of complex labels.
\end{proof}

We make explicit one other algebraic point used below. Define a null
raw label to be one whose insertion gives zero correlation with every
finite list of exterior insertions at distinct points. A slab operator
containing a null label is zero: pair it with two arbitrary word vectors,
use the null-label property on the resulting whole correlation, and use
density. Consequently \(M(w)\) depends only on the quotient classes of
its labels. Assume that the orthogonal action on those classes is an
actual group action and agrees, for signed coordinate transformations,
with the raw signed-label action after projection. These are the precise
compatibilities used when composing rotations below.

\begin{proposition}[Orthogonal covariance in each character]
\label{ex-individual-rotations}
Under these compatibilities and Proposition~\ref{ex-ordinary-reflection},
\(M(Kw)=M(w)\) for every \(K\in O(4)\) and distinct-point word \(w\).
Thus \(S_\chi(Kw)=S_\chi(w)\) for every character.
\end{proposition}
\begin{proof}
Write \(F(p,w)=\langle\Omega,M(p)M(w)\Omega\rangle\).
We first show simultaneous invariance
\begin{equation}\label{ex-pair-rotations}
 F(Kp,Kw)=F(p,w)\qquad(K\in O(4)).
\end{equation}
The cluster argument in Proposition~\ref{ex-ground-commutative} already
proves this whenever there exists a vector \(b\) with \(b_0>0\) and
\((Kb)_0>0\). Put \(a=(Ke_0)_0\). If \(a>-1\), choose
\(b=e_0+K^{-1}e_0\); orthogonality gives
\(b_0=(Kb)_0=1+a>0\). If \(a=-1\), then
\((\rho Ke_0)_0=1\), so the argument applies to \(\rho K\).
Individual invariance under \(\rho\) removes the extra reflection
from each factor and again gives \eqref{ex-pair-rotations}.

For any \(J\in O(4)\), rotate a pair backwards by \(J^{-1}\),
apply ordinary reflection to its second factor alone, and rotate forwards
by \(J\). Equation~\eqref{ex-pair-rotations} gives
\[
 F(p,w)=F(J^{-1}p,J^{-1}w)
       =F(J^{-1}p,\rho J^{-1}w)
       =F(p,J\rho J^{-1}w).
\]
There is no ambiguity from raw representatives, by quotient independence.
Pair values separate compressed operators: if \(F(p,w)=F(p,z)\) for
every \(p\), replace \(p\) by \(\theta v\) to obtain
\[
 \langle M(v)\Omega,(M(w)-M(z))\Omega\rangle=0.
\]
The linear span of \(M(v)\Omega\) is dense in \(\mathcal H_0\),
so \((M(w)-M(z))\Omega=0\); separation gives \(M(w)=M(z)\).
Thus \(M\) is individually invariant under every conjugate
\(J\rho J^{-1}\).

These conjugates are all hyperplane reflections: map \(e_0\) to any
chosen unit normal by an orthogonal map. For completeness, hyperplane
reflections generate \(O(d)\). Given an orthogonal map \(K\), if
\(Ke_1\ne e_1\), reflection in the hyperplane normal to
\(Ke_1-e_1\) sends \(Ke_1\) to \(e_1\); this follows immediately
from the reflection formula and \(|Ke_1|=|e_1|=1\). Composing with
this reflection gives a map fixing \(e_1\). It preserves
\(e_1^\perp\), and induction on dimension writes its restriction as
a product of reflections there, extended to fix \(e_1\). If it already
fixes \(e_1\), start with that restriction. Dimension zero is trivial.
Invariance under products follows from the group action on quotient
labels. This proves individual invariance under every orthogonal map,
and evaluation by a character gives the asserted phase covariance.
\end{proof}

Ordinary conjugation now follows as well. Since \(\theta w\) is a
permutation of \(\rho\overline w\),
\[
 S_\chi(\overline w)=S_\chi(\rho\overline w)
 =S_\chi(\theta w)=\overline{S_\chi(w)}.
\]
Signed coordinate covariance follows from the compatibility of the
quotient action with the raw signed action. Transporting time-axis
reflection positivity by signed permutations gives reflection positivity
about every coordinate axis: choose a signed permutation taking that
axis to the time axis, rotate each positive word and its labels, and use
covariance and the conjugation compatibility in the Gram matrix.
Lemma~\ref{ex-phase-smeared-rp} then supplies the corresponding smeared
positivity. The unproved quantitative input in this chapter is the
bounded real approximation property on arbitrarily fine generic grids;
its uniformity over finitely many replicated copies is indispensable.
